\section{Cycle intervals, parity, and orbit bookkeeping}
\label{app:cycles}

The affine coordinates in \cref{thm:cycle-solver} make positivity
transparent.  Under even compatibility,
\[
  n_i=(-1)^{i-1}x+b_i.
\]
Every odd index contributes a strict lower bound $x>-b_i$, and every even
index contributes a strict upper bound $x<b_i$.  Both index sets are
nonempty for even $r\ge2$, so the interval $(L(q),U(q))$ is finite.  It may
be empty.  Since all $q_i$ and hence all $b_i$ are integers, its positive
solution count is exactly
\[
  \#(\Z\cap(L(q),U(q)))
  =\max\{0,U(q)-L(q)-1\}.
\]
For the odd block, the recurrence modulo two reads
$n_{i+1}\equiv n_i\pmod2$, because $q_i$ is even.  Thus filtering $x=n_1$
to odd values filters every vertex simultaneously.

At $r=1$, \cref{eq:odd-candidate} gives $n_1=q_1/2$.  This is exactly the
loop at a power of two.  In the odd block only $q_1=2$ remains, corresponding
to the loop at vertex $1$.  This length-one convention is independent of the
restriction $r\ge2$ used for Hilbert--Schmidt trace powers.

\paragraph{Based words, rotations, and primitive period.}
A solution of \cref{eq:cycle-system} is a based vertex word
$(n_1,\ldots,n_r)$.  Rotating the word changes its base point but represents
the same cyclic orbit.  The edge labels rotate with it.  A word is primitive
only when its least cyclic shift period is $r$; a repeated traversal is not a
new primitive orbit.  The label tuple is derived from the vertex word and
does not replace it: the same compatible even label tuple can have several
vertex solutions, as $(4,8,8,4)$ demonstrates.

The coefficient $1/r$ in the local determinant logarithm removes cyclic base
points at the level of based trace terms.  It does not declare every label
tuple primitive, nor does it turn a repeated word into a new primitive.
Any Euler-product reorganization would require a separate least-period
reduction.  No such primitive-orbit product is claimed in this paper.

\paragraph{Finite-cutoff evaluation.}
For a cutoff $N$, a matrix-independent evaluator lists the allowed ordered
labels, applies \cref{thm:cycle-solver}, and then retains precisely the
solutions with $1\le n_i\le N$ for every $i$.  Summing
$\prod_i n_i^{-s}$ over these based solutions agrees with the trace of the
cutoff matrix power.  The procedure is exact because the theorem classifies
all positive solutions for each fixed ordered tuple; it uses neither a
spectral approximation nor a primitive-orbit heuristic.
